% xctest-and-swift-testing.tex — the two Apple test frameworks side by side: the XCTestCase
% lifecycle, XCTAssert → #expect / #require, traits, arguments, exit tests, attachments,
% parallelism, XCTest-only performance tests, the leak test, and what Swift 6.0 → 6.4 added.
% Sources (checked 2026-10-09):
%   developer.apple.com/documentation/xctest/set-up-and-tear-down-state-in-your-tests (class setUp
%     once; setUp() async throws → setUpWithError() → setUp(); test; teardown blocks LIFO; tearDown()
%     → tearDownWithError() → tearDown() async throws; class tearDown; assertions allowed in instance
%     setUp / tearDown, not in the class methods; blocks run even with continueAfterFailure = false;
%     nothing runs after a crash; async variants on an arbitrary actor unless @MainActor)
%   developer.apple.com/documentation/xctest/defining-test-cases-and-test-methods (instance method,
%     no parameters, no return value, name begins with "test") · /xctest/xctestcase (async tests) ·
%     /xctestcase/continueafterfailure (default true) · /xctestcase/executiontimeallowance (600 s,
%     rounded up to a minute, timeouts must be enabled) · /xctexpectedfailure/options/isstrict
%     (default true) · /xctattachment/lifetime (default deleteOnSuccess) · /xctmetric (7 metric
%     classes) · /xctmeasureoptions/iterationcount (iterationCount + 1 runs, first ignored)
%   developer.apple.com/library/archive …/testing_with_xcode chapters 01 + 04 (XCTest with Xcode 5;
%     performance tests vs a baseline Xcode 6; UI tests + recording Xcode 7; baseline = average +
%     max standard deviation, stored per device configuration; its reading list dates them: WWDC
%     2013/409 Xcode 5, 2014/414 Xcode 6, 2015/406 Xcode 7)
%   github.com/swiftlang/swift-book TSPL.docc/LanguageGuide/AutomaticReferenceCounting.md ("ARC
%     automatically sets a weak reference to nil when the instance … is deallocated" — the leak test)
%   developer.apple.com/videos/play/wwdc2018/403 (Xcode 10 parallel distributed testing: simulator
%     clones, each test class on one runner) · wwdc2021/10296 (Xcode 13 test repetitions) ·
%     wwdc2024/10179 (keep XCTest for XCUIApplication, XCTMetric,
%     Objective-C; a separate suite instance per @Test; in-process parallelism, also on devices;
%     don't mix assertion APIs) · wwdc2024/10195 (max two argument collections; zip; parallel and
%     randomised order by default; .serialized inherited; XCTest = processes, one test each; @Tag)
%   developer.apple.com Xcode release notes: 13 (async test methods; async setUp first, async
%     tearDown last; three repetition modes) · 26 (Swift 6.2; exit tests; attachments in the test
%     report; XCTIssue warning severity) · 26.4 (Swift 6.3; image attachments; Issue severity; interop
%     off by default) · 27 (Swift 6.4; cross-framework failures shown as warnings, test-plan setting;
%     swift test --maximum-repetitions / --repeat-until)
%   swift.org/blog: announcing-swift-6 (2024-09-17; Swift Testing in the toolchain, all Swift
%     platforms, swift test runs both) · swift-6.1-released (2025-03-31, Xcode 16.3) ·
%     swift-6.2-released (2025-09-15; raw-identifier names) · swift-6.3-released (2026-03-24) ·
%     swift-6.4-released (2026-09-15)
%   github.com/swiftlang/swift-testing Sources/Testing/Testing.docc: MigratingFromXCTest (assertion
%     map; init / deinit, deinit only on class or actor, never async or throws; struct or actor
%     preferred; XCTest runs sync tests on the main actor, Swift Testing on an arbitrary task; interop
%     modes + default table) · OrganizingTests (@Suite optional; zero-argument init; distinct instance
%     per test; no @available on a suite; traits inherited) · DefiningTests · Parallelization
%     (--no-parallel; .serialized) · ParameterizedTesting (product of two collections; zip; selective
%     rerun needs CustomTestArgumentEncodable / RawRepresentable / Encodable / Identifiable) ·
%     known-issues · exit-testing · testing-for-errors-in-swift-code · AddingTags ·
%     developer.apple.com/documentation/testing confirmation(…) (checked when the closure returns)
%   github.com/swiftlang/swift-evolution proposals/testing: ST-0004 (time limit in minutes, shortest
%     wins; 6.0) · ST-0005 (ranged confirmations; 6.1) · ST-0006 (#expect(throws:) returns the
%     error; 6.1) · ST-0007 (TestScoping; 6.1) · ST-0008 (exit tests; 6.2; not iOS) · ST-0009
%     (attachments; 6.2) · ST-0011 (issue handling traits; 6.2) · ST-0012 (exit-test captures,
%     Sendable + Codable, explicit list; 6.3) · ST-0013 (severity; 6.3) · ST-0014 (images; 6.3) ·
%     ST-0016 (Test.cancel; 6.3) · ST-0021 (interop; 6.4) · ST-0022 (CustomTestReflectable; 6.4) ·
%     ST-0023 (Transferable attachments; 6.4) · ST-0024 (per-case repetitions; 6.4) · ST-0025
%     (the existing swift test --filter / --skip options; their tag: prefix is 6.5, unreleased —
%     left out) · github.com/swiftlang/swift-package-manager Documentation (--filter / --skip
%     regular expressions)
% @labels: area=testing kind=api platform=apple topic=testing discussion=267
% @tags: xctest, swift-testing, xctestcase, addteardownblock, expect-macro, require-macro, parameterized-tests, test-traits, exit-tests, test-attachments, confirmation, parallel-testing, xctmetric, retain-cycle
\documentclass[8pt]{extarticle}
\usepackage{printup-sheet}
\usepackage{array}

\newcommand\ct[1]{\texttt{#1}}
\lstdefinelanguage{SwiftSheet}{
  morekeywords={protocol,class,final,struct,enum,func,var,let,weak,init,override,extension,
    if,else,return,guard,self,nil,try,await,async,throws,private,some,in,do,static,
    true,false,AnyObject,Void,String,Bool,Int,UInt,import},
  sensitive=true, morecomment=[l]{//}, morecomment=[s]{/*}{*/}, morestring=[b]"}

\tikzset{
  st/.style={box, font=\scriptsize, inner sep=1.2pt, minimum height=7mm, minimum width=11mm},
  once/.style={st, draw=sheetBrown, fill=sheetBrown!10},
  per/.style={st, draw=sheetBlue, fill=sheetBlue!8},
  tst/.style={st, draw=sheetGreen, fill=sheetGreen!14, font=\scriptsize\bfseries},
  tdn/.style={st, draw=sheetOrange, fill=sheetOrange!10},
  l/.style={font=\tiny, text=black!75, inner sep=1pt, align=center},
  hd/.style={font=\bfseries\scriptsize, text=sheetBlue, anchor=west, inner sep=0pt},
  tk/.style={circle, fill=#1, inner sep=0pt, minimum size=1.6mm},
  dot/.style={circle, draw=sheetBlue, inner sep=0pt, minimum size=2mm},
}

\begin{document}

\sheettitle{XCTest \& Swift Testing — lifecycle, macros, traits, what changed}{testing · folio}

\oneliner{\textbf{XCTest}: subclass \ct{XCTestCase}, name methods \ct{test…}, assert with
\ct{XCTAssert*}; a fresh instance per test method. \textbf{Swift Testing} (open source, in the
Swift toolchain, all Swift platforms): \ct{@Test} functions, any type is a suite, two macros
\ct{\#expect} / \ct{\#require}, traits and arguments — run \textbf{in parallel, in one process}
by default. Both live in one target; UI automation, \ct{XCTMetric} performance tests and
Objective-C stay in XCTest.}

\vspace{2pt}
\noindent\begin{tikzpicture}[sheet]
  % ===== XCTest row
  \node[hd] at (-0.05,0.95) {XCTest — one class, every test method};
  \node[once] (cs) at (0.6,0) {class\\\ct{setUp()}};
  \node[per] (a1) at (2.05,0) {\ct{setUp()}\\async throws};
  \node[per] (a2) at (3.4,0) {\ct{setUp}\\\ct{WithError()}};
  \node[per] (a3) at (4.6,0) {\ct{setUp()}};
  \node[tst] (t) at (5.8,0) {\ct{test…()}};
  \node[tdn] (b) at (7.1,0) {teardown\\blocks, LIFO};
  \node[tdn] (d1) at (8.45,0) {\ct{tearDown()}};
  \node[tdn] (d2) at (9.75,0) {\ct{tearDown}\\\ct{WithError()}};
  \node[tdn] (d3) at (11.15,0) {\ct{tearDown()}\\async throws};
  \node[once] (ct) at (12.6,0) {class\\\ct{tearDown()}};
  \foreach \a/\b in {cs/a1,a1/a2,a2/a3,a3/t,t/b,b/d1,d1/d2,d2/d3,d3/ct} \draw[flow] (\a) -- (\b);
  \node[draw=sheetGrey, dashed, rounded corners=2pt, inner sep=1.3mm, fit=(a1)(d3)] (per) {};
  \draw[flow, sheetBlue] (d3.north) -- ++(0,0.32) -| node[l, pos=0.25, above, text=sheetBlue]
     {next test method — on a \textbf{new} \ct{XCTestCase} instance} (a1.north);
  % ===== Swift Testing row
  \node[hd] at (-0.05,-0.8) {Swift Testing — one suite, every \ct{@Test}};
  \node[once] (s0) at (0.6,-1.5) {suite\\type};
  \node[per] (si) at (3.4,-1.5) {\ct{init()}\\async throws};
  \node[tst] (st) at (5.8,-1.5) {\ct{@Test func}};
  \node[tdn] (sd) at (9.1,-1.5) {\ct{deinit} — class\\or actor only};
  \draw[flow] (s0) -- (si); \draw[flow] (si) -- (st); \draw[flow] (st) -- (sd);
  \node[draw=sheetGreen!60!black, dashed, rounded corners=2pt, inner sep=1.3mm, fit=(si)(sd)] (sc) {};
  \node[l, anchor=north west, text=sheetGreen!45!black, align=left] at (sc.south west)
     {a \ct{TestScoping} trait (6.1) wraps the whole box — async set-up \emph{and} tear-down;
      \ct{deinit} can't be async or throw};
  \node[l, anchor=west, align=left] at (10.55,-1.5) {a distinct suite instance\\per \ct{@Test};
     struct, class or actor\\— struct preferred};
  % ===== notes
  \node[l, anchor=north west, align=left, text width=28mm] at (13.45,1.0)
     {\textcolor{sheetBrown}{brown} = once per class · \textcolor{sheetBlue}{blue} set-up ·
      \textcolor{sheetGreen!60!black}{green} the test · \textcolor{sheetOrange}{orange} tear-down.\\[1pt]
      Assertions count in instance set-up/tear-down; class methods have no instance — none there.\\[1pt]
      Tear-down runs after a failure, even with \ct{continueAfterFailure = false};
      \textcolor{sheetRed}{a crash skips all of it}.};
\end{tikzpicture}

\vspace{1pt}
\noindent\begin{tikzpicture}[sheet]
  \draw[thick, sheetBrown!60] (0,0) -- (5.55,0);
  \draw[thick, sheetBlue!60] (5.85,0) -- (16.5,0);
  \draw[sheetGrey] (5.6,-0.1) -- (5.7,0.1) (5.75,-0.1) -- (5.85,0.1);
  \foreach \x/\y/\col/\w/\t in {
     0.55/{2013\\Xcode 5}/sheetBrown/10.5mm/{XCTest},
     1.65/{2014\\Xcode 6}/sheetBrown/10.5mm/{\ct{measure} vs a baseline},
     2.75/{2015\\Xcode 7}/sheetBrown/10.5mm/{UI tests + recording},
     3.85/{2018\\Xcode 10}/sheetBrown/10.5mm/{parallel on simulator clones},
     4.95/{2021\\Xcode 13}/sheetBrown/10.5mm/{\ct{async} tests; repetition modes},
     6.75/{Sep 2024\\Swift 6.0 · Xcode 16}/sheetBlue/18mm/{\textbf{Swift Testing}: \ct{@Test}, \ct{\#expect}, traits, arguments, tags},
     8.7/{Mar 2025\\6.1 · Xcode 16.3}/sheetBlue/18mm/{\ct{TestScoping} traits; \ct{\#expect(throws:)} returns the error; ranged \ct{confirmation}},
     10.65/{Sep 2025\\6.2 · Xcode 26}/sheetBlue/18mm/{exit tests; attachments; raw-identifier names; issue-handling traits},
     12.6/{Mar 2026\\6.3 · Xcode 26.4}/sheetBlue/18mm/{\ct{severity: .warning}; \ct{Test.cancel}; image attachments; exit-test captures},
     14.55/{Sep 2026\\6.4 · Xcode 27}/sheetBlue/21.5mm/{XCTest ↔ Swift Testing interop; \ct{CustomTestReflectable}; per-case repetitions}} {
    \node[tk=\col] at (\x,0) {};
    \node[l, anchor=south, text=\col, font=\tiny\bfseries] at (\x,0.08) {\y};
    \node[l, anchor=north, text width=\w] at (\x,-0.08) {\t};
  }
\end{tikzpicture}

\begin{multicols}{2}
\raggedright
\footnotesize\setstretch{1.0}

\section{XCTest — the rules}
\begin{itemize}
  \item \textbf{Test method}: instance method of an \ct{XCTestCase} subclass, no parameters, no
        return value, name starts with lowercase \ct{test}; may be \ct{throws} or \ct{async}.
        Synchronous tests run on the \textbf{main actor}; \ct{async} set-up/tear-down and
        blocks run on an arbitrary actor unless marked \ct{@MainActor}.
  \item \textbf{Set-up}: \ct{setUpWithError()} when it can throw — a throw is recorded as a
        failure, no \ct{try!}; \ct{setUp() async throws} for async work.
        \ct{addTeardownBlock \{ \}} puts clean-up next to the code that needs it.
  \item \ct{continueAfterFailure} (default \ct{true}); \ct{false} ends the method at the first
        failure — the other tests still run. \ct{try XCTUnwrap(x)} fails and stops one test;
        \ct{x!} on \ct{nil} crashes the process: no tear-down.
  \item \textbf{Skip}: \ct{throw XCTSkip("…")} · \ct{try XCTSkipIf(c)} / \ct{…Unless(c)}.
        \textbf{Known failure}: \ct{XCTExpectFailure("…")} — \ct{isStrict} defaults to
        \ct{true}, so a failure that stops happening becomes an issue. Xcode 26: an
        \ct{XCTIssue} of warning severity does not fail the test.
  \item \ct{executionTimeAllowance}: default 600\,s, rounded \emph{up} to whole minutes, active
        only when timeouts are on (test plan or \ct{-test-timeouts-enabled YES}).
  \item \ct{XCTAttachment} + \ct{add(\_:)}: \ct{lifetime} defaults to \ct{.deleteOnSuccess} —
        kept only for failures unless \ct{.keepAlways}.
  \item \textbf{Parallel}: the simulator is cloned; each runner process takes \emph{whole
        classes}, one test at a time — one giant class never parallelises.
\end{itemize}

\section{Performance tests — XCTest only}
\ct{measure \{ \}} times the block; \ct{measure(metrics:options:)} records any of
\ct{XCTClockMetric} (time) · \ct{XCTCPUMetric} · \ct{XCTMemoryMetric} (physical) ·
\ct{XCTStorageMetric} (logical writes) · \ct{XCTOSSignpostMetric} (a signposted interval) ·
\ct{XCTApplicationLaunchMetric} · \ct{XCTHitchMetric}. The block runs
\ct{iterationCount}\,+\,1 times; the first is ignored (warm caches). A \textbf{baseline} =
average + allowed standard deviation, set from the report and stored \textbf{per device
configuration}.

\section{The leak test — weak, scope, nil}
\begin{lstlisting}[language=SwiftSheet]
func test_viewModel_isReleased() {
  weak var weakSUT: FeedViewModel?
  do {
    let sut = FeedViewModel(api: StubAPI())
    weakSUT = sut                 // observe, don't retain
    sut.start()                   // may build a cycle
  }                               // last strong ref gone
  XCTAssertNil(weakSUT)           // Swift Testing:
}                                 // #expect(weakSUT == nil)
func trackForLeaks(_ o: AnyObject, line: UInt = #line) {
  addTeardownBlock { [weak o] in XCTAssertNil(o, line: line) }
}                                 // after the test body
\end{lstlisting}

\section{Mixing the two (Swift 6.4, ST-0021)}
\ct{SWIFT\_TESTING\_XCTEST\_INTEROP\_MODE}: \ct{none} · \ct{limited} (an \ct{XCTAssert}
failing in a \ct{@Test} = warning) · \ct{complete} (= failure) · \ct{strict}
(\ct{fatalError}). Toolchain < 6.4: none — the failure was silently lost; packages with
tools-version ≥ 6.4 get complete, older ones limited; Xcode 27 sets it per test plan.

\section{Traps}
\begin{itemize}
  \trap{\ct{trackForLeaks} on a SUT kept in a \ct{var sut} property always fails: teardown
        blocks run \emph{before} \ct{tearDown()} nils it — build it in a local \ct{makeSUT()}.}
  \trap{Shared statics, singletons, \ct{UserDefaults.standard} race once tests run in
        parallel — inject them, or \ct{.serialized} as the last resort.}
  \trap{\ct{.timeLimit} / \ct{executionTimeAllowance} catch hangs, not slowness — that is
        \ct{measure}'s job.}
\end{itemize}

\columnbreak

\section{XCTest → Swift Testing}
{\scriptsize\setlength{\tabcolsep}{2pt}
\begin{tabular}{@{}>{\raggedright\arraybackslash}p{24mm}>{\raggedright\arraybackslash}p{50mm}@{}}
\toprule
\textbf{XCTest} & \textbf{Swift Testing}\\
\midrule
\ct{XCTAssertTrue(x)} & \ct{\#expect(x)}\\
\ct{XCTAssertEqual(a, b)}, \ct{…Nil}, \ct{…Identical}, \ct{…GreaterThan} & \ct{\#expect(a == b)}, \ct{x == nil}, \ct{===}, \ct{>} — one macro; a failure shows every operand's value\\
\ct{XCTAssertThrowsError} & \ct{let e = \#expect(throws: E.self) \{ try f() \}} — \ct{e: E?} (6.1); \ct{(any Error).self} for any\\
\ct{XCTAssertNoThrow} & \ct{\#expect(throws: Never.self) \{ try f() \}}\\
\ct{try XCTUnwrap(x)} & \ct{try \#require(x)} — throws \ct{ExpectationFailedError}, ends this test\\
\ct{XCTFail("…")} & \ct{Issue.record("…")}; \ct{severity: .warning} (6.3)\\
\ct{XCTSkipIf} & \ct{.enabled(if:)} · \ct{.disabled("why")}; mid-test \ct{try Test.cancel()} (6.3)\\
\ct{XCTExpectFailure} & \ct{withKnownIssue \{ \}} — no issue ⇒ records one, unless \ct{isIntermittent: true}\\
\ct{XCTestExpectation} & \ct{await confirmation(expectedCount: 1...) \{ c in \}} — counted when the closure \emph{returns}; it never waits\\
\ct{setUp} / \ct{tearDown} & \ct{init()} / \ct{deinit}; shared async → \ct{TestScoping} trait\\
\ct{executionTimeAllowance} & \ct{.timeLimit(.minutes(n))} — minutes only, shortest wins\\
\ct{XCTAttachment} & \ct{Attachment.record(value, named:)} (6.2)\\
\ct{measure}, \ct{XCUIApplication}, Objective-C & — stay in XCTest\\
\bottomrule
\end{tabular}}

\section{Swift Testing — the rules}
\begin{itemize}
  \item \ct{@Test} on a function at file scope or in a type; \ct{async}, \ct{throws},
        \ct{@MainActor} allowed. Name: \ct{@Test("…")} or a raw identifier
        \ct{func \textasciigrave adds two numbers\textasciigrave()} (6.2).
  \item A type holding tests \emph{is} a suite; \ct{@Suite} only adds a name or traits.
        Zero-argument \ct{init} required; never \ct{@available} on a suite. Suite traits
        (tags, \ct{.disabled}) are inherited.
  \item \ct{\#expect} expands to code capturing the source text and each subexpression's
        value; \ct{\#require} also throws. Mid-test failure → \ct{Issue.record}.
  \item \textbf{Parallel, randomised, in-process} (tasks), on devices too. Tests no longer
        default to the main actor. \ct{.serialized} — inherited by nested suites; on a
        parameterized test its cases run one by one.
  \item \textbf{Run}: \ct{swift test} runs both frameworks; \ct{--filter} / \ct{--skip} take a
        regex; \ct{--no-parallel}. 6.4: \ct{--maximum-repetitions n --repeat-until pass|fail}
        repeats only the matching cases — as Xcode 27's repetition mode now does.
  \item \textbf{Traits}: \ct{.enabled(if:)} · \ct{.disabled} · \ct{.bug("url")} ·
        \ct{.tags(.slow)} with \ct{extension Tag \{ @Tag static var slow: Self \}} (test
        plans include/exclude by tag) · \ct{.timeLimit} · \ct{.serialized} ·
        \ct{.compactMapIssues} / \ct{.filterIssues} (6.2) · a custom \ct{TestScoping} (6.1).
  \item \textbf{Exit tests} (6.2): \ct{await \#expect(processExitsWith: .failure) \{ … \}} —
        the closure is \ct{main()} of a child process of the same test binary;
        \ct{.success} · \ct{.failure} · \ct{.exitCode(n)} · \ct{.signal(n)}; may observe
        stdout/stderr. Captures (6.3): listed, \ct{Sendable} + \ct{Codable}. macOS, Linux,
        FreeBSD, OpenBSD, Windows — \textbf{not iOS}.
  \item \textbf{Attachments}: \ct{String}, bytes, \ct{Data}, \ct{Encodable} (6.2); \ct{CGImage},
        \ct{UIImage}, \ct{NSImage}, \ct{CIImage} as \ct{.png}/\ct{.jpeg} (6.3); \ct{Transferable}
        (6.4). Xcode: test report + \ct{.xcresult}; \ct{swift test --attachments-path}.
\end{itemize}

\section{Arguments — one test case each}
\noindent\begin{tikzpicture}[sheet]
  \foreach \ox/\zip/\ttl in {0/0/{\ct{arguments: [1,2,3], [a,b,c]}}, 3.05/1/{\ct{arguments: zip(…)}}} {
    \node[l, anchor=west] at (\ox,1.35) {\ttl};
    \foreach \i/\a in {0/1,1/2,2/3} \node[l] at (\ox+0.55+\i*0.4,1.05) {\a};
    \foreach \j/\b in {0/a,1/b,2/c} \node[l] at (\ox+0.2,0.75-\j*0.3) {\b};
    \foreach \i in {0,1,2} \foreach \j in {0,1,2} {
      \pgfmathparse{(\zip==0 || \i==\j) ? 1 : 0}
      \ifnum\pgfmathresult=1 \node[dot, fill=sheetBlue!45] at (\ox+0.55+\i*0.4,0.75-\j*0.3) {};
      \else \node[dot, draw=sheetGrey!50] at (\ox+0.55+\i*0.4,0.75-\j*0.3) {}; \fi }
  }
  \node[l, anchor=west, text=sheetBlue] at (1.75,0.45) {9 cases};
  \node[l, anchor=west, text=sheetBlue] at (4.8,0.45) {3 cases};
\end{tikzpicture}\par
Each element is its own case, run in parallel, shown and re-runnable alone (if the argument is
\ct{CustomTestArgumentEncodable}, \ct{RawRepresentable}, \ct{Encodable} or \ct{Identifiable}).
At most \textbf{two} collections — every combination; \ct{zip} pairs them.

\end{multicols}

\noindent{\footnotesize\color{sheetGrey}\textit{Related:} Testing fundamentals · Test doubles ·
Testable design · Testing ViewModels, Coordinators, UI \& snapshots · Async \& network testing ·
Flaky tests, CI, coverage \& TDD · ARC, strong / weak / unowned · Instruments · Macros \& result
builders · Sanitizers \& runtime checkers}

\end{document}
